\documentclass{sstt-2026-abstract}
%% The class already loads standard commonly-used packages:
%% amsmath, hyperref, biblatex, tikz-cd, etc.
%%
%% To help us preparing a book of abstracts, please avoid loading further
%% packages unless absolutely essential.

\addbibresource{sstt-example.bib}

%%%%%%%%%%%%%%%%
%% Extra commands, theorem types (\newtheorem), etc
%% Please keep simple and use sparingly

%%%%%%%%%%%%%%%%
%% Abstract metadata
%% These replace the usual \title, \author commands

\abstracttitle{Some theorems on type theories}

\abstractspeaker{Alice Speaker} % speaker
\abstractcoauthor{Bob Coauthor} % coauthor(s)/collaborator(s) as needed
% \abstractcoauthor{Charlie Third}

\begin{document}

%%%%%%%%%%%%%%%%
%% Abstract main text

Abstract text to follow here. Abstracts should be concise: maximum 1 page, including any references.

In this talk we will demonstrate:

\begin{theorem*}
  Every type theory is normalising, canonical, and conservative.
\end{theorem*}

We will make use of techniques including algebra, coalgebra, rewriting, gluing, reflection, refraction, and continuity principles \cite{martin-lof:bibliopolis,cartmell:generalised-algebraic-theories}

\lipsum[3-4]

%%%%%%%%%%%%%%%%
%% Bibliography, if needed:
\printbibliography

\end{document}